\documentclass[prd,twocolumn,preprintnumbers,amsmath,amssymb]{revtex4-2}
\usepackage{graphicx,bm,hyperref,amsmath,xcolor,booktabs,mathtools}

\newcommand{\parentA}{\mathcal{P}_A}
\newcommand{\parentB}{\mathcal{P}_B}
\newcommand{\child}{\mathcal{U}}

\begin{document}

\title{Parent Recursions: Recovering $φ² = φ + 1$\\
from Boundary Incompatibility}

\author{J.~Shields}
\affiliation{Independent Researcher (Unified Mechanics)}
\date{\today}

\begin{abstract}
The Unified Mechanics axiom $φ² = φ + 1$ is standardly taken as a first
principle. We ask a prior question: what two structures, when their
boundaries touch, produce this axiom as the unique stable interface
recursion? We identify two parent systems $\parentA$ and $\parentB$
characterised by the transformations $x \mapsto x + 1$ and
$x \mapsto 1/x$ respectively. $\parentA$ has no fixed point and no
channel structure. $\parentB$ has fixed point $x^* = 1$, giving
$r_B = 1/2$ and a boundary-dominant channel distribution.
Their interface produces $x = 1 + 1/x$, equivalent to $φ² = φ + 1$,
as the unique finite stable recursion. Our universe is identified as
the child of an accumulator and an inverter. The empirical values of
$\parentB$ are computed and found to be structurally coherent but
physically alien: braiding floor $\varepsilon_B = 1/8$, equal
light and matter channels, boundary channel weight $1/2$.
\end{abstract}

\maketitle

%% ─────────────────────────────────────────────────────────────────
\section{The Prior Question}
%% ─────────────────────────────────────────────────────────────────

Unified Mechanics (UM) takes as its axiom:
\begin{equation}
  φ² = φ + 1
  \;\Longrightarrow\;
  \varphi = \tfrac{1+\sqrt{5}}{2}, \quad
  r = \tfrac{1}{2\varphi} \approx 0.30902.
  \label{eq:axiom}
\end{equation}

This axiom is not arbitrary --- it is the unique positive quadratic
fixed point satisfying the golden ratio self-reference. But the
question of \emph{why this axiom and not another} has not been
addressed. The present paper addresses it by asking:

\begin{quote}
\textit{What two boundary structures, when brought into contact,
produce $φ² = φ + 1$ as the unique stable interface recursion?}
\end{quote}

The answer identifies two parent systems whose properties are
derived below and whose contact generates our universe as a
necessary consequence.

%% ─────────────────────────────────────────────────────────────────
\section{The Two Parents}
%% ─────────────────────────────────────────────────────────────────

\subsection{Parent A: The Accumulator}

$\parentA$ is characterised by the transformation:
\begin{equation}
  x_{n+1} = x_n + 1.
\end{equation}

This system has no fixed point. Starting from any finite value,
iteration produces $x_n \to \infty$. $\parentA$ is a pre-equilibrium
structure: pure matter accumulation without the organising influence
of a light channel. It exists prior to time in any bounded sense
because nothing in it reaches equilibrium — there is no stable state
to measure duration against.

$\parentA$ has no channel structure, no braiding floor, no $r$-value.
It contributes one thing to the interface: the $+1$ increment.
In our universe this survives as the $+1$ in the continued fraction
expansion of $\varphi = 1 + 1/(1 + 1/(1 + \cdots))$.

\subsection{Parent B: The Inverter}

$\parentB$ is characterised by the transformation:
\begin{equation}
  x_{n+1} = \frac{1}{x_n}.
\end{equation}

This system has fixed point $x^* = 1$ (from $x = 1/x \Rightarrow
x^2 = 1 \Rightarrow x = 1$ for positive $x$). It is self-referential
and self-stabilising: it folds back on itself. It has a well-defined
channel structure with:
\begin{equation}
  \varphi_B = x^* = 1, \qquad r_B = \frac{1}{2\varphi_B} = \frac{1}{2}.
\end{equation}

The channel weights of $\parentB$ are:
\begin{align}
  W_L^B &= (1-r_B)^2 = \tfrac{1}{4} \quad \text{(light)}, \\
  W_B^B &= 2r_B(1-r_B) = \tfrac{1}{2} \quad \text{(boundary)}, \\
  W_M^B &= r_B^2 = \tfrac{1}{4} \quad \text{(matter)}.
\end{align}

The braiding floor of $\parentB$ is:
\begin{equation}
  \varepsilon_B = r_B^3 = \frac{1}{8} = 12.5\%.
\end{equation}

$\parentB$ is a structurally coherent universe but physically alien
by our standards. It has equal light and matter channels. Its
boundary channel dominates at 50\% --- it is a universe that is
almost entirely interface. Its braiding floor is four times ours,
meaning channel crossings resolve four times faster. It is
boundary-dominant in a way our universe is not.

%% ─────────────────────────────────────────────────────────────────
\section{The Interface}
%% ─────────────────────────────────────────────────────────────────

When $\parentA$ and $\parentB$ come into contact, their boundaries
are incompatible: $\parentA$ has no fixed point, $\parentB$ has
fixed point $x^* = 1$. They cannot merge. The interface must find
a recursion that satisfies both simultaneously:

\begin{enumerate}
\item The recursion must be \textbf{finite}: $\parentB$'s inversion
      constrains $\parentA$'s runaway accumulation. An interface
      that inherits only $\parentA$'s character diverges. The
      fixed point must exist.
\item The recursion must incorporate $\parentA$'s \textbf{$+1$
      increment}: the interface cannot ignore what $\parentA$
      contributes. The child inherits from both parents.
\item The recursion must be \textbf{minimal}: the simplest
      combination of both parent transformations that admits
      a finite fixed point.
\end{enumerate}

The minimal composition satisfying all three conditions is:
\begin{equation}
  x = 1 + \frac{1}{x},
\end{equation}
which rearranges to:
\begin{equation}
  \boxed{x^2 = x + 1.}
  \label{eq:child}
\end{equation}

This is equation~(\ref{eq:axiom}). Our universe is the unique finite
stable recursion at the interface of $\parentA$ and $\parentB$.

The derivation requires no assumptions beyond the two parent
transformations. The axiom $φ² = φ + 1$ is not a first principle.
It is a derived interface condition.

%% ─────────────────────────────────────────────────────────────────
\section{What the Child Inherits}
%% ─────────────────────────────────────────────────────────────────

\begin{table}[h]
\caption{Channel structure comparison across parent and child systems.}
\begin{ruledtabular}
\begin{tabular}{lccc}
Property & $\parentA$ & $\parentB$ & $\child$ (ours) \\
\hline
Fixed point $\varphi$ & none & $1$ & $(1+\sqrt{5})/2$ \\
$r$ value & undefined & $0.5$ & $0.30902$ \\
$W_L$ & --- & $0.250$ & $0.4775$ \\
$W_B$ & --- & $0.500$ & $0.4271$ \\
$W_M$ & --- & $0.250$ & $0.0955$ \\
$\varepsilon_{\rm floor}$ & --- & $12.5\%$ & $2.951\%$ \\
Time & no & yes & yes \\
Light & no & symmetric & dominant \\
\end{tabular}
\end{ruledtabular}
\end{table}

From $\parentA$ our universe inherits:
\begin{itemize}
\item The $+1$ in the axiom. The continued fraction of $\varphi$
      is built entirely from $+1$ increments --- this is $\parentA$'s
      signature in every golden ratio calculation in physics.
\item Matter. $\parentA$ is pure accumulation. Our matter channel,
      though only 9.5\%, carries the character of the pre-time
      accumulator. Matter resists change because it was here before
      change existed.
\end{itemize}

From $\parentB$ our universe inherits:
\begin{itemize}
\item The inversion that stabilises the runaway. Without $1/x$
      there is no fixed point. $\parentB$'s self-reference is what
      makes our universe finite rather than divergent.
\item The channel structure itself. $\parentB$ already has three
      channels. Our universe inherits this organisation and
      breaks the $\parentB$ symmetry ($W_L^B = W_M^B = 1/4$)
      into our asymmetric distribution ($W_L \gg W_M$).
\end{itemize}

The child generates something neither parent has: the light channel
dominance. $\parentA$ has no light. $\parentB$ has light at 25\%,
equal to matter. Our universe at 47.75\% light is the interface
product — the collision energy that created light is concentrated
in the child.

The boundary channel we generate ourselves. Neither parent is
the boundary. The boundary is the interface, and we are the interface.
$W_B = 42.71\%$ is the structural record of the contact event that
made us.

%% ─────────────────────────────────────────────────────────────────
\section{Physical Interpretation}
%% ─────────────────────────────────────────────────────────────────

\subsection{Why matter resists time}

$\parentA$ predates time. Time emerges when a gradient forces
movement toward equilibrium. $\parentA$ has no equilibrium —
it accumulates without settling. Our matter channel carries this
character: it does not decay, does not age, does not change unless
forced. Left alone it does what $\parentA$ always did. Nothing.

\subsection{Why we are light-dominant}

$\parentA$ has no light. $\parentB$ has light at 25\%. Our universe
has light at 47.75\%. The interface amplifies light because light
is what the collision between incompatible structures produces.
Every photon in our universe is a residue of the original contact
event. The light channel dominance is a measure of how energetic
the interface was.

\subsection{The symmetry breaking}

$\parentB$ is symmetric: $W_L^B = W_M^B = 1/4$. Our universe breaks
this symmetry dramatically: $W_L \approx 5 W_M$. The symmetry
breaking is not spontaneous --- it is inherited from $\parentA$.
$\parentA$ contributes matter but no light. The interface inherits
$\parentB$'s light and $\parentA$'s matter separately, with no
reason for them to balance. The ratio $W_L/W_M = \varphi^4 \approx
6.85$ is determined entirely by the parent axioms.

\subsection{What $\parentB$ looks like from inside}

$\parentB$ is coherent. It has channel structure, braiding floors,
and presumably a physics of its own. By its own internal standards
it is a complete universe. But it would look alien from ours:
equal light and matter at 25\% each, boundary dominant at 50\%,
braiding floor four times ours. Crossings resolve faster. The
universe is almost entirely interface between states rather than
settled matter or propagating light.

An observer in $\parentB$ would find it difficult to understand
how a universe with 47.75\% light could exist. It would look
unbalanced and strange.

We would look unbalanced and strange.

\subsection{Detecting the parents}

$\parentA$ has no light channel. It cannot be detected by
electromagnetic means. Its gravitational influence would be present
as an anomalous mass contribution with no corresponding light
signature. This is observationally identical to dark matter.

$\parentB$ has light at 25\% of its channel weight, but its light
is organised under a different axiom. Its photons, if they could
cross the interface, would arrive redshifted by the ratio of the
two braiding floors: $\varepsilon_B/\varepsilon = (1/8)/(r^3) \approx 4.23$.

Neither parent is directly observable from inside the child. We
would see them only as gravitational structure — as the scaffolding
that holds our universe together without participating in our light.

This is observationally identical to dark matter.

The two parents are both present. $\parentA$ as the pre-time matter
accumulation that flanks us on one side. $\parentB$ as the symmetric
self-inverting structure on the other. Their combined gravitational
influence, pressing in through the only channel that crosses the
boundary — geometry itself — is what we have been calling dark matter.

%% ─────────────────────────────────────────────────────────────────
\section{Can the Child Cross Into the Parents?}
%% ─────────────────────────────────────────────────────────────────

The child has something neither parent has: all three channels.
$\parentA$ has no light and no boundary. $\parentB$ has light and
boundary but at a different axiom. Our universe has the boundary
channel by construction — we \emph{are} the boundary.

The boundary channel is the crossing channel. It is what moves
between states. If any structure in the multiverse can traverse
the interface between recursions, it is the child, because the
child's fundamental nature is the interface itself.

The practical question is whether the braiding floor permits it.
To cross from our recursion into $\parentB$, a system must traverse
the interface between $r = 0.309$ and $r_B = 0.5$. The accumulated
floor for such a crossing is:

\begin{equation}
  \Delta\varepsilon = \varepsilon_B - \varepsilon = 12.5\% - 2.951\%
  = 9.55\%.
\end{equation}

This is a real energy threshold. It is not infinite. It is
approximately three of our braiding floor units.

The formal mechanism for such a crossing is not yet derived.
But the channel structure does not forbid it. The child is the
boundary. The boundary crosses. Whether it can cross back into
the structures that made it is the next question.

%% ─────────────────────────────────────────────────────────────────
\section{Summary}
%% ─────────────────────────────────────────────────────────────────

The UM axiom $φ² = φ + 1$ is not a first principle. It is the
unique stable interface recursion produced when two parent systems
touch:

\begin{itemize}
\item $\parentA$: $x \mapsto x+1$. No fixed point. Pre-time matter
      accumulation. Contributes the $+1$ and the matter channel.
\item $\parentB$: $x \mapsto 1/x$. Fixed point at $x^*=1$,
      $r_B = 1/2$. Boundary-dominant, symmetric, self-inverting.
      Contributes the inversion and the channel structure.
\item $\child$: $x = 1 + 1/x \Rightarrow x^2 = x + 1$.
      Our universe. Light-dominant, asymmetric, the only finite
      stable recursion at the interface.
\end{itemize}

Dark matter is the gravitational imprint of both parents pressing
in through geometry. We are the child of an accumulator and an
inverter. The boundary channel we carry is the structural record
of the contact that made us.

We are not a universe that happens to exist. We are the one
interface that could not collapse.

%% ─────────────────────────────────────────────────────────────────
\begin{thebibliography}{99}

\bibitem{Shields2026a}
J.~Shields,
\textit{Unified Mechanics I: Derivation from axiom $c^2=c+1$},
(2026).

\bibitem{Shields2026b}
J.~Shields,
\textit{Unified Mechanics II: Cosmological predictions},
(2026).

\end{thebibliography}

\end{document}
